\section{Experiments}
We benchmark three models against the primary temporal split: the
hand-written multiplicative-update NMF ($k=16$ latent factors), a shrunk
cosine item-$k$NN baseline, and a non-personalised popularity baseline.
Ranking quality is measured at $K=10$ with Recall@$K$, NDCG@$K$, Mean
Average Precision at $K$ (MAP@$K$), and catalogue coverage: the fraction of
recipes that appear in at least one user's top-$K$ list across the
evaluated population. Candidates for every model are the full recipe
catalogue minus the user's training history.

To measure the nutrition filter's effect in isolation from any ranking
quality difference between models, we fix the model to NMF, request a
candidate pool of 50 recipes per user, apply the filter
(\texttt{max\_calories=450}, \texttt{max\_sodium\_mg=600}, reflecting a
"light meal, moderate sodium" profile), and compare NDCG@10 and Recall@10
on the first 10 survivors against the unfiltered first 10 candidates from
the same pool. We additionally report the mean fraction of the 50-recipe
candidate pool that the filter removes, which is the direct cost of
enforcing the constraint.

All numbers in Section~\ref{sec:results} are produced by
\texttt{scripts/export\_results\_table.py} against the synthetic fixture
described in \texttt{scripts/download\_data.py}, run with no network access
or Kaggle credentials, so the pipeline, this paper, and the CI workflow all
execute identically.
